% Curriculum vitae for Pei Wu.
% This file is maintained directly in LaTeX; LyX is not required.
\documentclass[10pt,american]{article}

% Page and text setup.
\usepackage[T1]{fontenc}
\usepackage[utf8]{inputenc}
\usepackage[letterpaper,left=0.78in,right=0.78in,top=0.72in,bottom=0.68in]{geometry}
\usepackage{amsmath}
\usepackage{amssymb}
\usepackage{babel}
\usepackage{fancyhdr}
\usepackage{lmodern}
\usepackage{tabularx}
\usepackage{xcolor}
\usepackage{hyperref}

\frenchspacing
\setlength{\parindent}{0pt}
\setlength{\parskip}{0pt}
\renewcommand{\arraystretch}{1.15}

% Document-wide values.
\newcommand{\cvname}{Pei Wu}
\newcommand{\cvyear}{2026}
\newcommand{\ac}{\mathrm{AC}}
\newcommand{\me}{\textbf{P. Wu}}
\definecolor{cvaccent}{RGB}{184,0,55}

\hypersetup{
  colorlinks=true,
  urlcolor=cvaccent,
  pdftitle={Pei Wu - Curriculum Vitae},
  pdfauthor={Pei Wu},
  pdfsubject={Academic curriculum vitae},
  pdfkeywords={Pei Wu, theoretical computer science, complexity theory}
}
\urlstyle{same}

% Running header and footer. The first page uses only the page number.
\pagestyle{fancy}
\fancyhf{}
\fancyhead[C]{\small\itshape \cvname\ \textbullet\ Curriculum Vitae\ \textbullet\ \cvyear}
\fancyfoot[C]{\thepage}
\renewcommand{\headrulewidth}{0.4pt}
\renewcommand{\footrulewidth}{0pt}
\setlength{\headheight}{14pt}

\fancypagestyle{firstpage}{%
  \fancyhf{}
  \fancyfoot[C]{\thepage}
  \renewcommand{\headrulewidth}{0pt}
}

% Reusable CV layout primitives.
\newcolumntype{D}[1]{>{\raggedright\arraybackslash}p{#1}}
\newcolumntype{Y}{>{\raggedright\arraybackslash}X}

\newcommand{\cvsection}[1]{%
  \par\addvspace{1.25em}%
  {\large\bfseries\color{cvaccent}#1}\par\nobreak
  \vspace{2pt}\hrule height 0.8pt\nobreak\vspace{0.65em}\nobreak
}

\newcommand{\cventry}[3]{%
  \begin{tabularx}{\linewidth}{@{}D{0.16\linewidth}Y@{}}
    #1 & \textsc{#2}\\[-1pt]
       & #3
  \end{tabularx}\par\vspace{0.35em}%
}

\newcommand{\cvpublication}[4]{%
  \begin{tabularx}{\linewidth}{@{}D{0.052\linewidth}Y@{}}
    \textbf{[#1]} & \textbf{\textit{#2}}\\[-1pt]
                   & #3\\[-1pt]
                   & #4
  \end{tabularx}\par\vspace{0.45em}%
}

\newcommand{\cvwork}[3]{%
  \begin{minipage}{\linewidth}
    \raggedright
    \textbf{\textit{#1}}\\[-1pt]
    #2%
    \def\cvdetail{#3}%
    \ifx\cvdetail\empty\else\\[-1pt]#3\fi
  \end{minipage}\par\vspace{0.75em}%
}

\newcommand{\cvtalk}[3]{%
  \begin{tabularx}{\linewidth}{@{}D{0.16\linewidth}Y@{}}
    #1 & \textbf{\textit{\lq\lq#2\rq\rq}}\\[-1pt]
       & #3
  \end{tabularx}\par\vspace{0.35em}%
}

\newcommand{\cvperson}[2]{%
  \textsc{#1}\\[-1pt]
  #2\par\vspace{0.4em}%
}

\begin{document}
\raggedright
\thispagestyle{firstpage}

\begin{center}
  {\LARGE \cvname}\\[2pt]
  {\Large Assistant Professor}\\[5pt]
  {\large\textsc{Department of Computer Science and Engineering}}\\[1pt]
  {\large\textsc{Pennsylvania State University}}\\[7pt]
  \href{mailto:pei.wu@psu.edu}{pei.wu@psu.edu}
  \quad\textbullet\quad
  \href{https://pwu.netlify.app/}{pwu.netlify.app}
  \quad\textbullet\quad
  \href{https://scholar.google.co.uk/citations?user=i152hTQAAAAJ}{Google Scholar}
\end{center}

\cvsection{Academic Appointments}

\cventry{2024--present}{Penn State University, University Park, Pennsylvania}{%
  \textit{Assistant Professor}}

\cventry{2023--2024}{Weizmann Institute of Science, Rehovot, Israel}{%
  Faculty of Mathematics and Computer Science\newline
  \textit{Senior Postdoctoral Fellow; supervisor: Thomas Vidick}}

\cventry{Summer 2023}{Simons Institute for the Theory of Computing, Berkeley, California}{%
  \textit{Research Fellow}}

\cventry{2021--2023}{Institute for Advanced Study, Princeton, New Jersey}{%
  School of Mathematics\newline
  \textit{Postdoctoral Member; supervisor: Avi Wigderson}}

\cvsection{Education}

\cventry{2015--2021}{University of California, Los Angeles}{%
  \textit{Ph.D., Computer Science}\newline
  Thesis: \textit{Communication and Computation}; advisor: Alexander Sherstov}

\cventry{2013--2015}{Dartmouth College}{%
  \textit{M.S., Computer Science}; thesis advisor: Amit Chakrabarti}

\cventry{2009--2013}{Nanjing University, China}{%
  \textit{B.S., Computer Science and Technology}}

\cvsection{Publications\hfill{\normalfont\footnotesize\itshape\color{black}Author order is alphabetical.}}

\cvpublication{11}{Randomized and quantum lifting for one-way conservative NOF model}%
  {H. Wang, \me}%
  {The 51st International Symposium on Mathematical Foundations of Computer Science (MFCS 2026)}

\cvpublication{10}{Quantum algorithms on edge lists: hiding, shuffling, and cycle finding}%
  {A. S. Gilani, D. Wang, \me, X. Zhou}%
  {The 53rd International Colloquium on Automata, Languages, and Programming (ICALP 2026). \textit{Also accepted as a contributed talk at TQC 2026}}

\cvpublication{9}{Quantum Merlin--Arthur with an internally separable proof}%
  {R. Bassirian, B. Fefferman, I. Leigh, K. Marwaha, \me}%
  {The 21st Conference on the Theory of Quantum Computation, Communication and Cryptography (TQC 2026), contributed talk}

\cvpublication{8}{Coherence in property testing: quantum-classical collapses and separations}%
  {F. G. Jeronimo, N. Magrafta, J. Slote, \me}%
  {The 28th Annual Quantum Information Processing Conference (QIP 2025)}

\cvpublication{7}{Dimension-independent disentanglers from unentanglement and applications}%
  {F. G. Jeronimo, \me}%
  {The 39th Computational Complexity Conference (CCC 2024)}

\cvpublication{6}{Pseudorandom and pseudoentangled states from subset states}%
  {F. G. Jeronimo, N. Magrafta, \me}%
  {The 19th Conference on the Theory of Quantum Computation, Communication and Cryptography (TQC 2024), poster}

\cvpublication{5}{An optimal \lq\lq it ain't over till it's over\rq\rq\ theorem}%
  {R. Eldan, A. Wigderson, \me}%
  {The 55th ACM Symposium on Theory of Computing (STOC 2023)}

\cvpublication{4}{The power of unentangled quantum proofs with non-negative amplitudes}%
  {F. G. Jeronimo, \me}%
  {The 55th ACM Symposium on Theory of Computing (STOC 2023)}

\cvpublication{3}{An optimal separation of randomized and quantum query complexity}%
  {A. A. Sherstov, A. A. Storozhenko, \me}%
  {The 53rd ACM Symposium on Theory of Computing (STOC 2021). \textit{SIAM Journal on Computing, 52(2):525--567, 2023}}

\cvpublication{2}{Near-optimal lower bounds on the threshold degree and sign-rank of $\ac^{0}$}%
  {A. A. Sherstov, \me}%
  {The 51st ACM Symposium on Theory of Computing (STOC 2019). \textit{Invited to appear in SIAM Journal on Computing (special issue for STOC 2019)}}

\cvpublication{1}{Optimal interactive coding for insertions, deletions, and substitutions}%
  {A. A. Sherstov, \me}%
  {The 58th Annual Symposium on Foundations of Computer Science (FOCS 2017). \textit{IEEE Transactions on Information Theory, 65(10):5971--6000, 2019}}

\cvsection{Survey and Expository Article}

\cvwork{The QMA(2) universe: complexity, entanglement, and optimization}%
  {F. G. Jeronimo, I. Leigh, \me}%
  {\textit{Invited survey, ACM SIGACT News, 57(1):64--99, 2026}}

\cvsection{Manuscripts and Preprints}

\cvwork{Optimal quantum de Finetti theorems via argmax rounding}%
  {F. G. Jeronimo, \me, H. Xu}{}

\cvwork{An argmax principle for sum-of-squares relaxations on the sphere}%
  {F. G. Jeronimo, \me, H. Xu}{}

\cvwork{An optimal analysis of the product test}%
  {J. Beckey, F. G. Jeronimo, \me}{}

\cvwork{The power of unentanglement without destructive interference}%
  {Y. Liu, \me}{}

\cvsection{Advising}

\cvperson{Haoyu Wang}{Ph.D. student (in progress), Penn State University}
\cvperson{Haochen Xu}{Ph.D. student (in progress), Penn State University}
\cvperson{John Brenton}{Undergraduate researcher, Penn State University}

\newpage
\cvsection{Teaching}

\cventry{Fall 2026}{CMPSC 360}{Discrete Mathematics}
\cventry{Spring 2026}{CMPSC 564}{Complexity of Combinatorial Problems}
\cventry{Fall 2025}{CMPSC 464}{Theory of Computation}

\cvsection{Selected Talks and Seminars}

\cvtalk{Aug--Sep 2024}{Coherence in property testing: quantum-classical separations and collapses}{%
  Center on Frontiers of Computing Studies, Peking University; Theory Seminar, Penn State University}

\cvtalk{Mar 2024}{Multiple Merlins in the modern era}{%
  Theory Lunch, Weizmann Institute of Science}

\cvtalk{Apr 2023--Jan 2024}{The power of unentangled proofs with non-negative amplitudes}{%
  Simons Institute; Weizmann Institute of Science; The University of Texas at Austin; Nanjing University; QIP 2024}

\cvtalk{Jun 2023}{It ain't over till it's over}{%
  \lq\lq Analysis and TCS,\rq\rq\ Simons Institute}

\cvtalk{Sep--Nov 2022}{Random restrictions on Boolean functions with small influences}{%
  Princeton University; Shandong University; Nanjing University; DIMACS and Rutgers University}

\cvtalk{Nov 2022}{Polynomial method in communication complexity}{%
  CS/DM Seminar, Institute for Advanced Study}

\cvtalk{Oct 2021}{Recent progress on query complexity}{%
  Two lectures, CS/DM Seminar, Institute for Advanced Study}

\cvtalk{Apr 2021}{Optimal separation of randomized and quantum query complexity}{%
  Algorithm and Complexity Seminar, University of Waterloo (online)}

\cvtalk{Feb 2020}{Settling the threshold degree and sign-rank of $\ac^{0}$}{%
  Invited plenary talk, Southern California Theory Day, University of California, Riverside}

\cvsection{University Service}

\cventry{2025--2026}{Department Hiring Committee}{}
\cventry{2025}{Computer Science and Engineering Seminar Series}{Organizer}
\cventry{2025}{Graduate Qualifying Committee}{}

\cvsection{Professional Service}

\textbf{Program committee:} FSTTCS 2026

\vspace{0.35em}
\textbf{Conference reviewing:} STOC/FOCS, CCC, ITCS, ICALP, STACS, QIP, TQC

\vspace{0.35em}
\textbf{Journal referee:} PRX, SICOMP, TIT, \textit{Quantum}, \textit{Algorithmica}, JCSS

\end{document}
